\section{Schubert Cycles}

\bdefn

    Consider a complex flag $\V$ in $V$, i.e.\ a sequence of subspaces
    $$ 0 < V_1 < \cdots < V_{n-1} < V_n = V,\qquad \dim V_i = i . $$
    Furthermore suppose we have a sequence $a$
    $$ n-k \geq a_1 \geq \cdots \geq a_k \geq 0 . $$
    Define the {\emphcolor Schubert variety} to be the closed subset
    $$ \Sigma_a(\V) = \set{\Lambda\in G(k,n)}[\dim(V_{n-k+i-a_i}\cap\Lambda)\geq i\hbox{ for all $i$}] . $$

\edefn

\bremark

    In general for a subspace $\Gamma\leq V$, the set
    $$ \set{\Lambda\in G(k,n)}[\dim(\Gamma\cap\Lambda)\geq i] $$
    is closed, as $\dim(\Gamma\cap\Lambda)=\dim\Gamma+\dim\Lambda-\dim(\Gamma+\Lambda)$ so the condition
    is equivalent to $\dim(\Gamma+\Lambda)\leq\dim\Gamma+k-i$.
    Choosing bases for $\Lambda,\Gamma$, this is iff the rank of the matrix obtained by concatenating the
    bases has rank $\leq\dim\Gamma+k-i$, which is iff all submatrices of larger size have zero determinant.

    So we see that this set above is an algebraic subset of $G(k,n)$.
    And thus $\Sigma_a(\V)\subseteq G(k,n)$ is algebraic as well.
    We will show later that it is irreducible and thus a subvariety.

\eremark

For convenience, we write
$$ \Sigma_{a_1,\dots,a_s}(\V) = \Sigma_{a_1,\dots,a_s,0,\dots,0} $$
and
$$ \Sigma_{a^s} = \Sigma_{a,\dots,a} $$
to duplicate $a$ $s$ times.

\bexam

    Note that for any flag $\V$,
    $$ \Sigma_0(\V) = G(k,n) . $$
    That is, any $k$-dimensional subspace $\Lambda$ satisfies $\dim(V_{n-k+i}\cap\Lambda)\geq i$ for all $i$.
    Indeed,
    $$ \dim(V_{n-k+i}\cap\Lambda) = \dim V_{n-k+i} + \dim\Lambda - \dim(V_{n-k+i}+\Lambda) \geq
    n - k + i + k - n = i . $$

\eexam

\bexam

    Suppose $V_\ell$ is a subspace of dimension $\ell$, then extend it to an arbitrary complete flag $\V$
    and we see that
    $$ \Sigma_{n-k+1-\ell}(\V) = \set{\Lambda}[\Lambda\cap V_\ell\neq0] . $$
    Indeed for $i=1$ the Schubert condition gives
    $$ \dim(V_\ell\cap\Lambda) \geq 1 \iff V_\ell\cap\Lambda\neq 0 . $$
    The conditions on $i\geq1$ are tautological as shown by the previous example.

\eexam

\bexam[name={schubertinclusion}]

    Given an $\ell$-subspace $V_\ell$ then we have (for any extension $\V$)
    $$ \Sigma_{(n-\ell)^k}(\V) = \set{\Lambda}[\Lambda\subseteq V_\ell] , $$
    because the Schubert conditions gives for all $i$
    $$ \dim(\Lambda\cap V_{\ell-k+i}) \geq i , $$
    so for $i=k$ this means $\Lambda\subseteq V_\ell$, and then considering the first example the rest of
    the conditions follow.

    And similarly
    $$ \Sigma_{(n-k)^\ell}(\V) = \set{\Lambda}[V_\ell\subseteq\Lambda] $$
    because
    $$ \dim(\Lambda\cap V_i) \geq i \hbox{ for $i\leq\ell$} \iff V_\ell \leq \Lambda $$
    and the other conditions are tautological again.

    Thus we have that $\Sigma_{(n-k)^k}(\V)$ is the point $V_k$.

\eexam

Note that if we define for sequences $a\leq b$ iff $a_i\leq b_i$ for all $i$, then
$$ \Sigma_a(\V) \subseteq \Sigma_b(\V) \iff a \leq b . $$
For a sequence $a$ define
$$ \abs a = \sum_i a_i . $$

\bthrm

    The Schubert variety $\Sigma_a(\V)$ is an irreducible subvariety of $G(k,V)$ of codimension $\abs a$.

\ethrm

\bprop[name={intwithlin}]

    Let $e_1,\dots,e_n$ be a basis for $V$ and consider the standard flag $V_i=\gen{e_1,\dots,e_i}$.
    Let $a$ be a decreasing sequence and $b$ be the strictly increasing sequence $b_i=n-k+i-a_i$.
    Then $\Sigma_a(\V)\subseteq G(k,n)\subseteq\bP^N$ is the intersection
    $$ \Sigma_a(\V) = G(k,V) \cap \bigcap_{I\not\leq b}V(p_I) , $$
    where $p_I$ is the Pl\"ucker coordinate associated to $I=(i_1<\cdots<i_k)$.

\eprop

\bproof

    For $\Lambda\in G(k,V)$ write $\Lambda$ as the row space as the matrix $A=(a_{ij})$.
    The Pl\"ucker coordinate $p_I(\Lambda)$ is then the determinant of the submatrix of $A$ whose columns
    are given by $I$.

    Now suppose that $\Lambda\in\Sigma_a(V)$ so that $\dim(\Lambda\cap V_{b_i})\geq i$ for all $i$.
    For $I\not\leq b$ there is a $j$ such that $i_j>b_j$, and since $\dim(\Lambda\cap V_{b_j})\geq j$,
    there are $j$ linearly independent vectors in $\Lambda$ whose coordinates past $b_j$ are zero.
    If we take these vectors to be $v_1,\dots,v_j$ and extend them to a basis $v_1,\dots,v_k$ of $\Lambda$,
    then for every index in $I$ past $i_j$ (inclusive), the coordinates of $v_1,\dots,v_j$ are zero.
    Thus the submatrix of $A$ determined by $I$ has a $j\times(k-j+1)$ block of zeroes.
    That is, it is of the form

    \medskip
    $$ \left(\vcenter{\halign{\strut\hfil$#$\hfil&\@gobble{#1}\kern.2cm\vrule\kern.2cm&\hfil$#$\hfil&\rlap{\kern.5cm$#$}\cr
        \omit\hfil\rmht{\ $j-1$\ }\hfil&\omit&\omit\hfil\rmht{\ $k-j+1$\ }\hfil\cr
        *&&0&j\cr\noalign{\hrule}
        *&&*&k-j\cr
    }}\right)\kern1cm . $$
    Thus its rank is at most $(k-j)+(j-1)=k-1$ which means it is singular, so $p_I(\Lambda)=0$.

    Now suppose that $p_I(\Lambda)=0$ for all $I\not\leq b$, but $\dim(\Lambda\cap V_{b_j})<j$ for some $j$.
    So if consider the submatrix of $A$ determined by columns of index $>b_j$, it must have rank $>k-j$.
    This is because there are $<j$ rows which are also in $\Lambda\cap V_{b_j}$:

    \medskip
    $$ \left(\vcenter{\halign{\strut\hfil$#$\hfil&\@gobble{#1}\kern.2cm\vrule\kern.2cm&\hfil$#$\hfil&\rlap{\kern.5cm$#$}\cr
        \omit\hfil\rmht{\ $b_j$\ }\hfil\cr
        *&&\ 0\ &<j\cr\noalign{\hrule}
        *&&*&>k-j\cr
    }}\right)\kern1cm . $$

    Performing row operations gives us a block-diagonal matrix of full rank, and so the right column of blocks
    is linearly independent.
    So there are at least $k-j+1$ columns with indices greater than $b_j$ which are linearly independent.
    Extend these to $k$ linearly independent columns (because $A$ is of rank $k$) and let $I$ be the
    indices of these columns.
    Then the submatrix determined by $I$ is nonsingular so $p_I(\Lambda)\neq0$.
    Because at least $k-j+1$ of these indices is greater than $b_j$, and as they are ascending this means
    that $i_j>b_j$, a contradiction.
    \qed

\eproof

The reason for wanting to consider Schubert varieties follows from the above examples.
Consider the following problem.

\bprob[name={interfourq}]

    If $V_1,V_2,V_3,V_4\subseteq\bP^{2n+1}$ are four general $n$-planes, how many lines $L\subseteq\bP^{2n+1}$
    intersect all four?

\eprob

We know by the above example that $\Sigma_n(V_i)\subseteq\bG(1,2n+1)$ is the number of lines meeting $V_i$,
so we want to count the intersection
$$ Q = \bigcap\Sigma_n(V_i) . $$
Unfortunately we do not yet have the tools to do this.
But this motivates the following desired structure.
We want a ring $A$ which is generated by the subvarieties of $G(k,n)$, with a suitable equivalence so that
$\Sigma_a(\V)\in A$ depends only on $a$ and not on the choice of flag $\V$.
Furthermore the ring product should be given by the intersection of varieties.
Then in $A$ we have that $Q=\Sigma_n^4$ is a sum of points, and counting those points will give us the answer.

We now describe the construction of such a ring.

%Now consider \refmath[proposition]{intwithlin}: we have $b=(n+1,2n+2)$ which means the only $I$s which are
%not bound by $b$ are of the form $I=(i_1,i_2)$ with $i_1>n+1$.
%So we choose the indices from $(n+1,2n+2]$ which gives us $\binom{n+1}2$ choices for $I$.
%Thus $\Sigma_n(V_i)=\bG(1,2n+1)\cap W_i$ where $W_i$ is an $\parens{N-\binom{n+1}2}$-plane of $\bP^N$.
%Here we have
%$$ N = \binom{2n+2}2 - 1 = 2n^2 + 3n,\qquad \binom{n+1}2 = \frac{n^2+n}2 . $$
%Now the number of 

